% NGÂN HÀNG TUYỂN SINH 10 – BÀI 2

\begin{bt}%[TS10 2.0 2.1]%[ma:dc26679d]
Cho phương trình: $2x^2 - x - 5 = 0$.
\textbf{a)} Chứng tỏ phương trình trên có hai nghiệm phân biệt $x_1, x_2$. \par
\textbf{b)} Không giải phương trình, hãy tính giá trị của biểu thức: $A = x_1^2 + x_2^2 - x_1 - x_2$.
\loigiai{
\textbf{a)} Chứng tỏ phương trình có hai nghiệm phân biệt \hfill \textit{(0,25 điểm)} \par
Phương trình $2x^2 - x - 5 = 0$ có hệ số $a = 2$, $c = -5$.
Ta có: $a \cdot c = 2 \cdot (-5) = -10 < 0$.
Do đó, phương trình luôn có hai nghiệm phân biệt $x_1, x_2$. \par
\textbf{b)} Tính giá trị của biểu thức $A$ \hfill \textit{(0,75 điểm)} \par
Theo định lí Viète, ta có:
$\heva{&S = x_1 + x_2 = \frac{1}{2} \\ &P = x_1 x_2 = -\frac{5}{2}}$ \par
Biến đổi biểu thức $A$:
$A = (x_1^2 + x_2^2) - (x_1 + x_2)$
$A = (x_1 + x_2)^2 - 2x_1 x_2 - (x_1 + x_2)$ \par
Thay các giá trị $S$ và $P$ vào biểu thức:
$A = \left(\frac{1}{2}\right)^2 - 2\left(-\frac{5}{2}\right) - \frac{1}{2}$
$A = \frac{1}{4} + 5 - \frac{1}{2} = \frac{1}{4} + \frac{20}{4} - \frac{2}{4} = \frac{19}{4}$. \par
Vậy $A = \frac{19}{4}$.
}
\end{bt}

\begin{bt}%[TS10 2.0 2.1]%[ma:c72a9adb]
Cho phương trình: $3x^2 + 4x - 2 = 0$.
\textbf{a)} Chứng minh phương trình luôn có hai nghiệm phân biệt $x_1, x_2$. \par
\textbf{b)} Không giải phương trình, hãy tính giá trị của biểu thức: $B = \frac{x_1}{x_2 + 1} + \frac{x_2}{x_1 + 1}$.
\loigiai{
\textbf{a)} Chứng minh phương trình có hai nghiệm phân biệt \hfill \textit{(0,25 điểm)} \par
Phương trình $3x^2 + 4x - 2 = 0$ có $a = 3$, $c = -2$.
Ta thấy $a \cdot c = 3 \cdot (-2) = -6 < 0$.
Suy ra phương trình luôn có hai nghiệm phân biệt $x_1, x_2$. \par
\textbf{b)} Tính giá trị của biểu thức $B$ \hfill \textit{(0,75 điểm)} \par
Theo định lí Viète, ta có:
$x_1 + x_2 = -\frac{4}{3}$ và $x_1 x_2 = -\frac{2}{3}$. \par
Quy đồng mẫu thức biểu thức $B$:
$B = \frac{x_1(x_1 + 1) + x_2(x_2 + 1)}{(x_1 + 1)(x_2 + 1)}$
$B = \frac{x_1^2 + x_1 + x_2^2 + x_2}{x_1 x_2 + x_1 + x_2 + 1}$
$B = \frac{(x_1^2 + x_2^2) + (x_1 + x_2)}{x_1 x_2 + (x_1 + x_2) + 1}$
$B = \frac{(x_1 + x_2)^2 - 2x_1 x_2 + (x_1 + x_2)}{x_1 x_2 + (x_1 + x_2) + 1}$ \par
Thay số vào biểu thức:
Tử số: $\left(-\frac{4}{3}\right)^2 - 2\left(-\frac{2}{3}\right) + \left(-\frac{4}{3}\right) = \frac{16}{9} + \frac{4}{3} - \frac{4}{3} = \frac{16}{9}$.
Mẫu số: $-\frac{2}{3} + \left(-\frac{4}{3}\right) + 1 = -2 + 1 = -1$. \par
Vậy $B = \frac{16/9}{-1} = -\frac{16}{9}$.
}
\end{bt}

\begin{bt}%[TS10 2.0 2.1]%[ma:aabf4bbd]
Cho phương trình $x^2 - 5x - 3 = 0$.
\textbf{a)} Chứng tỏ phương trình trên có hai nghiệm phân biệt $x_1, x_2$. \par
\textbf{b)} Không giải phương trình, tính giá trị của biểu thức: $C = x_1^3 + x_2^3 - x_1 x_2$.
\loigiai{
\textbf{a)} Chứng tỏ phương trình có hai nghiệm phân biệt \hfill \textit{(0,25 điểm)} \par
Phương trình $x^2 - 5x - 3 = 0$ có $a = 1$, $c = -3$.
Vì $a \cdot c = 1 \cdot (-3) = -3 < 0$ nên phương trình luôn có hai nghiệm phân biệt $x_1, x_2$. \par
\textbf{b)} Tính giá trị của biểu thức $C$ \hfill \textit{(0,75 điểm)} \par
Áp dụng định lí Viète, ta có:
$x_1 + x_2 = 5$ và $x_1 x_2 = -3$. \par
Biến đổi biểu thức $C$:
$C = (x_1^3 + x_2^3) - x_1 x_2$
$C = (x_1 + x_2)^3 - 3x_1 x_2(x_1 + x_2) - x_1 x_2$ \par
Thay các giá trị vào ta được:
$C = (5)^3 - 3(-3)(5) - (-3)$
$C = 125 + 45 + 3 = 173$. \par
Vậy $C = 173$.
}
\end{bt}

\begin{bt}%[TS10 2.0 2.2]%[ma:078c602f]
Cho phương trình: $2x^2 - 7x - 6 = 0$.
\textbf{a)} Chứng minh phương trình trên có hai nghiệm phân biệt $x_1, x_2$. \par
\textbf{b)} Không giải phương trình, hãy tính giá trị của biểu thức: $D = \frac{2x_1^2 - 6}{x_2} + \frac{2x_2^2 - 6}{x_1}$.
\loigiai{
\textbf{a)} Chứng minh phương trình có hai nghiệm phân biệt \hfill \textit{(0,25 điểm)} \par
Xét phương trình $2x^2 - 7x - 6 = 0$ có $a = 2$, $c = -6$.
Do $a \cdot c = 2 \cdot (-6) = -12 < 0$, phương trình có hai nghiệm phân biệt $x_1, x_2$. \par
\textbf{b)} Tính giá trị của biểu thức $D$ \hfill \textit{(0,75 điểm)} \par
Theo định lí Viète: $x_1 + x_2 = \frac{7}{2}$ và $x_1 x_2 = -3$. \par
Vì $x_1, x_2$ là nghiệm của phương trình nên:
$2x_1^2 - 7x_1 - 6 = 0 \Rightarrow 2x_1^2 - 6 = 7x_1$.
Tương tự, $2x_2^2 - 6 = 7x_2$. \par
Thay vào biểu thức $D$, ta được:
$D = \frac{7x_1}{x_2} + \frac{7x_2}{x_1} = 7 \left(\frac{x_1^2 + x_2^2}{x_1 x_2}\right)$
$D = 7 \cdot \frac{(x_1 + x_2)^2 - 2x_1 x_2}{x_1 x_2}$ \par
Thay các giá trị tổng và tích vào:
$D = 7 \cdot \frac{\left(\frac{7}{2}\right)^2 - 2(-3)}{-3} = 7 \cdot \frac{\frac{49}{4} + 6}{-3} = 7 \cdot \frac{\frac{73}{4}}{-3}$
$D = 7 \cdot \left(-\frac{73}{12}\right) = -\frac{511}{12}$. \par
Vậy $D = -\frac{511}{12}$.
}
\end{bt}

\begin{bt}%[TS10 2.0 2.2]%[ma:b183a9c8]
Cho phương trình $3x^2 + 5x - 6 = 0$.
\textbf{a)} Chứng tỏ phương trình trên có hai nghiệm phân biệt $x_1, x_2$. \par
\textbf{b)} Không giải phương trình, hãy tính giá trị của biểu thức: $E = (3x_1^2 + 5x_1 - 7)^{2026} + (3x_2^2 + 5x_2 - 7)^{2026}$.
\loigiai{
\textbf{a)} Chứng tỏ phương trình có hai nghiệm phân biệt \hfill \textit{(0,25 điểm)} \par
Ta có $a = 3$, $c = -6$. Vì $a \cdot c = -18 < 0$ nên phương trình luôn có hai nghiệm phân biệt $x_1, x_2$. \par
\textbf{b)} Tính giá trị của biểu thức $E$ \hfill \textit{(0,75 điểm)} \par
Vì $x_1$ là nghiệm của phương trình $3x^2 + 5x - 6 = 0$ nên:
$3x_1^2 + 5x_1 - 6 = 0 \Rightarrow 3x_1^2 + 5x_1 = 6$.
Suy ra: $3x_1^2 + 5x_1 - 7 = 6 - 7 = -1$. \par
Tương tự đối với nghiệm $x_2$:
$3x_2^2 + 5x_2 - 6 = 0 \Rightarrow 3x_2^2 + 5x_2 = 6$.
Suy ra: $3x_2^2 + 5x_2 - 7 = 6 - 7 = -1$. \par
Thay vào biểu thức $E$, ta được:
$E = (-1)^{2026} + (-1)^{2026}$
Vì $2026$ là số mũ chẵn nên $(-1)^{2026} = 1$.
Vậy $E = 1 + 1 = 2$.
}
\end{bt}

\begin{bt}%[TS10 2.0 2.1]%[ma:32429cd2]
Cho phương trình: $x^2 + 3x - 5 = 0$.
\textbf{a)} Chứng minh phương trình trên có hai nghiệm phân biệt $x_1, x_2$. \par
\textbf{b)} Không giải phương trình, hãy tính giá trị của biểu thức: $F = \frac{x_1 - x_2}{x_1} \cdot \frac{x_2 - x_1}{x_2}$.
\loigiai{
\textbf{a)} Chứng minh phương trình có hai nghiệm phân biệt \hfill \textit{(0,25 điểm)} \par
Phương trình có $a = 1, c = -5 \Rightarrow a \cdot c = -5 < 0$.
Do đó, phương trình luôn có hai nghiệm phân biệt. \par
\textbf{b)} Tính giá trị của biểu thức $F$ \hfill \textit{(0,75 điểm)} \par
Theo định lí Viète: $x_1 + x_2 = -3$ và $x_1 x_2 = -5$. \par
Nhân phân phối biểu thức $F$:
$F = \frac{(x_1 - x_2)(x_2 - x_1)}{x_1 x_2} = \frac{-(x_1 - x_2)^2}{x_1 x_2}$
Mặt khác: $(x_1 - x_2)^2 = (x_1 + x_2)^2 - 4x_1 x_2$. \par
Thay số vào ta được:
$(x_1 - x_2)^2 = (-3)^2 - 4(-5) = 9 + 20 = 29$.
Do đó:
$F = \frac{-29}{-5} = \frac{29}{5}$. \par
Vậy $F = \frac{29}{5}$.
}
\end{bt}

\begin{bt}%[TS10 2.0 2.2]%[ma:a5b1abae]
Cho phương trình: $x^2 - 4x - 1 = 0$.
\textbf{a)} Chứng tỏ phương trình trên có hai nghiệm phân biệt $x_1, x_2$. \par
\textbf{b)} Không giải phương trình, hãy tính giá trị của biểu thức: $G = (x_1 - 2x_2)(x_2 - x_1) - x_1(x_1 + 2025x_2)$.
\loigiai{
\textbf{a)} Chứng tỏ phương trình có hai nghiệm phân biệt \hfill \textit{(0,25 điểm)} \par
Phương trình có $a = 1, c = -1 \Rightarrow a \cdot c = -1 < 0$.
Vậy phương trình có hai nghiệm phân biệt $x_1, x_2$. \par
\textbf{b)} Tính giá trị của biểu thức $G$ \hfill \textit{(0,75 điểm)} \par
Theo Viète, ta có $x_1 + x_2 = 4$ và $x_1 x_2 = -1$. \par
Thực hiện nhân phân phối và thu gọn biểu thức $G$:
$G = (x_1 x_2 - x_1^2 - 2x_2^2 + 2x_1 x_2) - (x_1^2 + 2025x_1 x_2)$
$G = 3x_1 x_2 - x_1^2 - 2x_2^2 - x_1^2 - 2025x_1 x_2$
$G = -2x_1^2 - 2x_2^2 - 2022x_1 x_2$
$G = -2(x_1^2 + x_2^2) - 2022x_1 x_2$
$G = -2[(x_1 + x_2)^2 - 2x_1 x_2] - 2022x_1 x_2$. \par
Thay các giá trị Viète vào:
$G = -2[4^2 - 2(-1)] - 2022(-1) = -2(16 + 2) + 2022$
$G = -2(18) + 2022 = -36 + 2022 = 1986$. \par
Vậy $G = 1986$.
}
\end{bt}

\begin{bt}%[TS10 2.0 2.1]%[ma:7569185f]
Cho phương trình $3x^2 - 5x - 4 = 0$.
\textbf{a)} Chứng minh phương trình trên có hai nghiệm phân biệt $x_1, x_2$. \par
\textbf{b)} Không giải phương trình, hãy tính giá trị của biểu thức: $H = x_1^2 + x_2^2 + x_1 x_2(2 + 2024x_1 + 2024x_2)$.
\loigiai{
\textbf{a)} Chứng minh phương trình có hai nghiệm phân biệt \hfill \textit{(0,25 điểm)} \par
Phương trình có hệ số $a = 3$, $c = -4$.
Vì $a \cdot c = 3 \cdot (-4) = -12 < 0$, phương trình có hai nghiệm phân biệt $x_1, x_2$. \par
\textbf{b)} Tính giá trị của biểu thức $H$ \hfill \textit{(0,75 điểm)} \par
Theo định lí Viète: $x_1 + x_2 = \frac{5}{3}$ và $x_1 x_2 = -\frac{4}{3}$. \par
Biến đổi biểu thức $H$:
$H = x_1^2 + x_2^2 + 2x_1 x_2 + 2024x_1 x_2(x_1 + x_2)$
$H = (x_1 + x_2)^2 + 2024x_1 x_2(x_1 + x_2)$ \par
Thay số vào biểu thức:
$H = \left(\frac{5}{3}\right)^2 + 2024 \left(-\frac{4}{3}\right) \left(\frac{5}{3}\right)$
$H = \frac{25}{9} - \frac{40480}{9} = -\frac{40455}{9}$. \par
Thu gọn phân số: $-\frac{40455}{9} = -13485$. \par
Vậy $H = -13485$.
}
\end{bt}

\begin{bt}%[TS10 2.0 2.1]%[ma:db1de1c6]
Cho phương trình: $x^2 - 2x - 7 = 0$.
\textbf{a)} Chứng minh phương trình có hai nghiệm phân biệt $x_1, x_2$. \par
\textbf{b)} Không giải phương trình, tính giá trị của biểu thức: $K = \frac{x_1^2}{x_2} + \frac{x_2^2}{x_1}$.
\loigiai{
\textbf{a)} Chứng minh phương trình có hai nghiệm phân biệt \hfill \textit{(0,25 điểm)} \par
Ta có $a = 1$, $c = -7$. Vì $a \cdot c = -7 < 0$, phương trình luôn có hai nghiệm phân biệt. \par
\textbf{b)} Tính giá trị của biểu thức $K$ \hfill \textit{(0,75 điểm)} \par
Theo định lí Viète: $x_1 + x_2 = 2$ và $x_1 x_2 = -7$. \par
Quy đồng mẫu thức biểu thức $K$:
$K = \frac{x_1^3 + x_2^3}{x_1 x_2}$
Mặt khác: $x_1^3 + x_2^3 = (x_1 + x_2)^3 - 3x_1 x_2(x_1 + x_2)$. \par
Thay số vào tử số:
$(x_1 + x_2)^3 - 3x_1 x_2(x_1 + x_2) = (2)^3 - 3(-7)(2) = 8 + 42 = 50$. \par
Thay vào biểu thức $K$:
$K = \frac{50}{-7} = -\frac{50}{7}$. \par
Vậy $K = -\frac{50}{7}$.
}
\end{bt}

\begin{bt}%[TS10 2.0 2.2]%[ma:d598b31a]
Cho phương trình $2x^2 + 5x - 6 = 0$.
\textbf{a)} Chứng tỏ phương trình trên có hai nghiệm phân biệt $x_1, x_2$. \par
\textbf{b)} Không giải phương trình, hãy tính giá trị của biểu thức: $M = \frac{x_1}{2x_2^2 + 5x_2 + 1} + \frac{x_2}{2x_1^2 + 5x_1 + 1}$.
\loigiai{
\textbf{a)} Chứng tỏ phương trình có hai nghiệm phân biệt \hfill \textit{(0,25 điểm)} \par
Ta thấy hệ số $a = 2$, $c = -6 \Rightarrow a \cdot c = -12 < 0$.
Do đó, phương trình có hai nghiệm phân biệt $x_1, x_2$. \par
\textbf{b)} Tính giá trị của biểu thức $M$ \hfill \textit{(0,75 điểm)} \par
Theo định lí Viète: $x_1 + x_2 = -\frac{5}{2}$ và $x_1 x_2 = -3$. \par
Vì $x_1, x_2$ là các nghiệm của phương trình nên:
$2x_1^2 + 5x_1 - 6 = 0 \Rightarrow 2x_1^2 + 5x_1 = 6 \Rightarrow 2x_1^2 + 5x_1 + 1 = 7$.
Tương tự: $2x_2^2 + 5x_2 + 1 = 7$. \par
Thay vào biểu thức $M$, ta được:
$M = \frac{x_1}{7} + \frac{x_2}{7} = \frac{x_1 + x_2}{7}$ \par
Thay tổng $x_1 + x_2 = -\frac{5}{2}$ vào biểu thức:
$M = \frac{-\frac{5}{2}}{7} = -\frac{5}{14}$. \par
Vậy $M = -\frac{5}{14}$.
}
\end{bt}
